\documentclass[11pt]{article}
\usepackage[margin=1in]{geometry}
\usepackage{amsmath,amssymb,amsthm}
\usepackage{booktabs}
\usepackage[hidelinks]{hyperref}

\newtheorem{theorem}{Theorem}
\newtheorem{lemma}{Lemma}
\newtheorem{corollary}{Corollary}
\theoremstyle{definition}
\newtheorem{definition}{Definition}

\newcommand{\Fp}{\mathbb{F}_p}
\newcommand{\Z}{\mathbb{Z}}
\newcommand{\Ehat}{\widehat{E}}
\newcommand{\red}{\mathrm{red}}
\newcommand{\ord}{\mathrm{ord}}
\newcommand{\Aut}{\mathrm{Aut}}
\DeclareMathOperator{\val}{v_p}
\DeclareMathOperator{\Uk}{U}

\title{\bf When a Control Does Not Control:\\
A Mechanism-Derived Refutation of a False Positive\\
in $p$-adic Cryptanalysis of the ECDLP}

\author{W.\ T.\ Trevena\thanks{Correspondence: william.todd.trevena@gmail.com.
Code and data: the \texttt{ecdsa-lift} repository.}}

\date{June 2026}

\begin{document}
\maketitle

\begin{abstract}
We report a systematic negative exploration of $p$-adic lifting attacks
on the elliptic curve discrete logarithm problem (ECDLP) for ordinary
curves with prime-order subgroups, of the kind underlying ECDSA. Across
forty experimental phases every structural probe of the canonical-lift /
formal-group attack surface came back inert. The methodological interest
lies in one episode. A Gowers $\Uk^2$ / Fourier analysis of the
Teichm\"uller lift error $\delta(k)$ produced a persistent
$5$--$11\sigma$ anomaly that a standard permutation control certified as
``real, index-ordered structure.'' We show that this conclusion was an
artifact of a control that cannot distinguish the hypotheses it was
meant to separate. Replacing it with two \emph{mechanism-derived} nulls
---one per candidate explanation---localizes the entire signal to a
single exact, universal, public identity (the antisymmetry
$\delta(k)+\delta(n-k)=[n]\tau(G)$), and shows it is neither new nor
cryptanalytically useful. We then prove, as an explicit instance of a
principle going back to Silverman and to Gadiyar--Padma, an inertness
theorem that explains why \emph{every} such anomaly on these curves must
reduce to artifact-or-identity: the only secret-dependent part of any
section lift error is governed by $k \bmod p^{e-1}$, which is
information-theoretically independent of the public residue $k \bmod n$
whenever $\gcd(n,p)=1$. The single inequality $\gcd(n,p)=1$ is sharp:
when it fails ($n=p$, anomalous curves) the obstruction vanishes and the
same machinery recovers $k$ with $100\%$ success---Smart's attack. The
contribution is less the (essentially folklore) theorem than the
discipline it certifies: in structure-hunting cryptanalysis, a control
is valid only against a named alternative hypothesis, and nulls must be
derived from mechanisms, not intuition.
\end{abstract}

\section{Introduction}

A recurring temptation in cryptanalysis is to lift a hard problem into a
richer setting---a $p$-adic ring, a formal group, a higher-genus
Jacobian---in the hope that structure invisible over $\Fp$ becomes
visible and exploitable. For the ECDLP this hope is well mapped:
Silverman's ``four faces of lifting'' \cite{silverman4faces,silvermanlift}
catalogues the local/global $\times$ torsion/non-torsion strategies and
the obstruction each meets, and the only $p$-adic local lift that ever
broke an ECDLP is Smart's attack on anomalous curves
\cite{smart,satoharaki,semaev}, where the group order equals the
characteristic.

This paper is a case study in \emph{how one knows when a lifting probe
has actually found nothing}. We conducted a long, deliberately
exhaustive negative exploration of the canonical-lift attack surface for
ordinary curves with large prime-order subgroups---the ECDSA regime---and
every phase was inert. Such a project invites a methodological hazard:
after dozens of null results, a probe eventually returns an
``anomaly,'' and the analyst must decide whether it is signal. We hit
exactly this case, initially mis-called it signal on the strength of a
textbook control, and only a more careful experiment overturned it. We
think the overturning, and the reason the anomaly \emph{had} to be
spurious, are more useful to report than the negative result itself.

\paragraph{Contributions.}
\begin{itemize}
\item A worked, quantitative example of a \emph{control that does not
control}: a permutation test that appears to certify structure but
provably cannot separate the competing explanations
(\S\ref{sec:falsepos}).
\item A general recipe---\emph{mechanism-derived controls}: enumerate the
candidate generating mechanisms and build one null per mechanism---and
its application, which localizes a $5$--$11\sigma$ Gowers anomaly to a
single exact public identity (\S\ref{sec:controls}).
\item An inertness theorem (\S\ref{sec:theorem}) that explains why such
anomalies are forced: for $\gcd(n,p)=1$, the secret-dependent part of any
section lift error lives in a coordinate ($k \bmod p^{e-1}$) independent
of the public one ($k \bmod n$). We state it as the explicit instance it
is of \cite{silverman4faces,gadiyarpadma}, with full proof and numerical
verification.
\item A sharp-boundary demonstration (\S\ref{sec:boundary}): the same
construction, run on an anomalous curve where $\gcd(n,p)=p$, recovers the
secret with $100\%$ success. The theorem's one hypothesis is exactly the
line between inert and broken.
\end{itemize}

We are explicit throughout that the theorem is not new mathematics; its
role is to certify, after the fact, that the methodological suspicion was
structurally warranted.

\section{Background and notation}\label{sec:background}

Let $p\ge 5$ be prime and $E/\Fp$ an elliptic curve with $\#E(\Fp)=N$.
Fix $G\in E(\Fp)$ of order $n$ with $\gcd(n,p)=1$, and $Q=kG$; the ECDLP
is to recover $k\in\Z/n\Z$ from $(E,G,Q)$. For $e\ge 1$, reduction modulo
$p$ on a fixed Weierstrass lift gives the exact sequence of abelian
groups
\begin{equation}\label{eq:ses}
0 \longrightarrow \Ehat(p\Z/p^e) \longrightarrow E(\Z/p^e)
  \xrightarrow{\ \red\ } E(\Fp) \longrightarrow 0,
\end{equation}
where the kernel of reduction $\Ehat(p\Z/p^e)$ is a $p$-group of order
$p^{e-1}$. The formal parameter $z=-X/Y$ identifies $\Ehat$ with
$(p\Z/p^e,+_F)$ ($+_F$ the formal group law), and the formal logarithm
$\log\colon\Ehat\to(p\Z/p^e,+)$ is an isomorphism with
$\log([m]u)=m\log(u)$. For $E\colon y^2=x^3+b$ the log has the sparse form
$\log(z)=z+\tfrac{b}{7}z^7+\cdots$, so for $z\in p\Z$ and $e\le 4$ one has
$\log(z)\equiv z\pmod{p^e}$: kernel arithmetic is additive in $z$ at our
precision.

A \emph{section} is a map $\tau\colon E(\Fp)\to E(\Z/p^e)$ with
$\red\circ\tau=\mathrm{id}$; the Teichm\"uller and Hensel lifts are
sections, generally not homomorphisms. The \emph{lift error} (a
$1$-cochain valued in $\Ehat$ by \eqref{eq:ses}) is
\begin{equation}
\delta_\tau(k)\ :=\ [k]\,\tau(G)\ -\ \tau(kG)\ \in\ \Ehat(p\Z/p^e).
\end{equation}
We write $\delta(k)$ for the Teichm\"uller case and
$z(\delta(k))\in p\Z/p^e$ for its formal parameter.

\paragraph{Prior art.} That a $p$-adic lift transports the ECDLP without
weakening it, that the canonical (Teichm\"uller) lift carries no
information, and that ``solving the DLP is equivalent to constructing a
non-canonical lift'' are due to Silverman \cite{silverman4faces,
silvermanlift} for elliptic curves and to Gadiyar--Maini--Padma and
Gadiyar--Padma \cite{gadiyarpadma} for $\Fp^\times$. Smart, Satoh--Araki
and Semaev \cite{smart,satoharaki,semaev} give the anomalous attack. The
canonical lift itself is Lubin--Serre--Tate \cite{lst}; its computational
instance is Satoh's point-counting algorithm \cite{satohcount}. Our
theorem is an explicit packaging of these in the lift-error language.

\section{The exploration, in brief}\label{sec:exploration}

The full study ran forty numbered phases against the canonical-lift /
formal-group attack surface for $j=0$ (CM by $\Z[\omega]$) and generic
ordinary curves. We summarize only the shape of the result; details are
in the repository. Probes included: injectivity and pseudorandomness of
$\delta$ (DFT flatness, maximal Berlekamp--Massey complexity, uniform
$p$-adic valuation, no polynomial/Mahler/multiplicative fit); a
homomorphic section from $H^2$ vanishing (which exists but, being a
homomorphism, annihilates the attack); an explicit Serre--Tate canonical
lift validated by $\delta_{\mathrm{can}}\equiv 0$; HNP/lattice (LLL,
BKZ) attacks on the cocycle; MOV embedding degree; ordinary/supersingular
bridges; genus-2 Jacobian covers; Weil pairings; Iwasawa $\mu$; and a
Frobenius/unit-root analysis. Every one was inert. Two exact relations
survived, both \emph{public} (computable from the curve) and both
giving no advantage: the automorphism-equivariance rank law
$\mathrm{rank}=\min(\cdot,(n-1)/|\Aut(E)|)$ for the leading digit of
$\delta$, and the antisymmetry
\begin{equation}\label{eq:antisym}
\delta(k)+\delta(n-k)=[n]\,\tau(G),
\end{equation}
verified to hold \emph{exactly} in the $z$-coordinate (e.g.\ for
$p=97$, $z(\delta(k))+z(\delta(n-k))$ takes a single value across all
$78$ pairs).

\section{A false positive: the Gowers anomaly}\label{sec:falsepos}

Phases of the study applied higher-order Fourier analysis to the digit
sequence of $\delta$. For digit $j$ encode
$f_j(k)=\exp\!\big(2\pi i\,\mathrm{digit}_j(z(\delta(k)))/p\big)$ and
measure the Gowers $\Uk^2$ norm against an $S^1$-random baseline. Across
$p=31$ to $521$ the $\Uk^2$ norm was elevated $5$--$11\sigma$; $\Uk^3$
and higher were at baseline. The leading Fourier coefficient decayed as
$|\hat f(\alpha)|\sim C n^{-\beta}$ with $\beta\in[0.28,0.47]<1/2$,
present on CM, $j{=}1728$, and non-CM curves alike. An internal report
concluded this was a genuine, previously unnoticed degree-$1$ structural
property of the Teichm\"uller cocycle ``not explained by any universal
identity.''

The evidence offered was a permutation control: randomly permuting the
\emph{values} of $f_j$ destroyed the elevation, which was read as proof
that the structure lived in the index ordering, hence was ``real.''
This inference is invalid, and seeing why is the point of the paper.

\section{Mechanism-derived controls}\label{sec:controls}

\paragraph{The principle.} A control is a null model, and a null model
tests exactly one alternative hypothesis: the one under which the
statistic is recomputed. The value-permutation control recomputes
$\Uk^2$ under ``the multiset of digit values, in random order.'' It
therefore separates ``structure in the value distribution'' from
``structure in the ordering''---and nothing else. But at least three
distinct mechanisms predict ordered structure here, and the permutation
kills all of them simultaneously, so it cannot attribute the signal to
any one:
\begin{enumerate}
\item[(H1)] \emph{Prefix-sum artifact.} $\delta(k)=-\sum_{i<k}c(iG,G)$ is
a running sum. Summation is a low-pass filter
($\hat s(\xi)=\hat d(\xi)/(1-e^{2\pi i\xi/n})$, a $1/\xi$ pole at
$\xi=0$), so \emph{any} increments with this multiset, once summed, carry
low-frequency mass and elevated $\Uk^2$.
\item[(H2)] \emph{Antisymmetry shadow.} The exact reflection
\eqref{eq:antisym} couples Fourier mode $\alpha$ with $-\alpha$, a rigid
constraint that elevates $\Uk^2$.
\item[(H3)] \emph{Genuine new structure}, the claimed explanation.
\end{enumerate}
The fix is to build one null per mechanism, each derived from how that
mechanism generates ordering, and to ask which null the real statistic
matches.

\paragraph{Control for H1 (increment shuffle).} $s(k)=z(\delta(k))$ is
the exact cumulative sum of its first differences $d(k)=s(k+1)-s(k)$.
Shuffle $d$, rebuild $s'$ by re-summing in the new order, recompute
$\Uk^2$; the increment multiset is identical, only the order changes.
Table~\ref{tab:p40} (Phase 40). For all primes with sequence length
$L\ge 66$ the shuffled-increment prefix sums sit \emph{at} the random
baseline ($z(\text{incShuf vs rand})$ within $\pm 1$, range $-0.26$ to
$+1.04$), while the true ordering sits $+4$ to $+12$ above them, growing
with $n$; the two smallest primes ($L\le 30$) are dominated by baseline
noise and excluded. So H1 is refuted: the digit nonlinearity destroys the
random walk's low-frequency mass, and the signal is in the \emph{specific}
order, not in summation.

\paragraph{Control for H2 (half-sequence).} The reflection
\eqref{eq:antisym} relates each $k$ to $n-k$. Restricting to the first
half $k=1,\dots,\lfloor(n-1)/2\rfloor$ removes every reflection pair from
the window, so H2 predicts no internal constraint there; H3 (an intrinsic
ordering structure) predicts the signal persists. Re-running the Phase 40
control on the half (Phase 40b): the elevation \emph{vanishes},
$z(\text{real vs incShuf})$ scattering in $[-2.4,+2.1]$ with no growth in
$n$. At matched sample size the contrast is decisive: $p=257$ gives
$z\approx+8$ on the full sequence ($L=256$) and $z\approx+0.1$ to $+0.6$
on the half ($L=128$).

\paragraph{Verdict.} The anomaly is the Fourier-domain shadow of the
exact, universal, public identity \eqref{eq:antisym}: H2, not H1 or H3.
It is real but not new, and---being a known public consequence of the
curve's geometry that yields only a $2$-to-$1$ search reduction---carries
no cryptanalytic advantage. The varying ``decay exponent''
$\beta\in[0.28,0.47]$ and its apparent curve-family dependence are
finite-size noise around a single reflection-induced spectrum; nothing
family-specific survives the half-sequence control.

\begin{table}[t]
\centering
\small
\begin{tabular}{rrrr}
\toprule
$p$ & $L$ & $z(\text{real vs incShuf})$ & $z(\text{incShuf vs rand})$\\
\midrule
67  & 78  & $+3.8$ to $+6.4$  & $-0.24$ to $+0.05$\\
97  & 78  & $+3.6$ to $+6.2$  & $+0.17$ to $+0.48$\\
167 & 166 & $+6.6$ to $+8.2$  & $+0.09$ to $+0.67$\\
257 & 256 & $+6.7$ to $+8.4$  & $-0.11$ to $+0.99$\\
521 & 520 & $+10.3$ to $+11.9$& $-0.07$ to $+0.35$\\
\bottomrule
\end{tabular}
\caption{Phase 40 increment-shuffle control (digits $j=1,2,3$). The
shuffled-increment prefix sums match the random baseline (right column
$\approx 0$); the true ordering is far above and grows with $n$ (middle).
H1 (prefix-sum artifact) is refuted. The half-sequence control
(Phase 40b) then collapses the middle column to noise, localizing the
signal to the antisymmetry \eqref{eq:antisym}.}
\label{tab:p40}
\end{table}

\section{Why it had to be a false positive: an inertness theorem}
\label{sec:theorem}

The episode above was not luck. For $\gcd(n,p)=1$ \emph{every}
secret-dependent feature of a section lift error is forced into a
coordinate the public data cannot see, so any apparent ``structure'' in
$\delta$ that is computable from $(E,G,Q)$ must reduce to a public
identity or an artifact. We make this precise.

\begin{lemma}[Canonical lift]\label{lem:can}
If $\gcd(n,p)=1$, the order-$n$ subgroup $\langle G\rangle$ lifts
uniquely to $\langle\widetilde G\rangle\subset E(\Z/p^e)$ of order $n$
with $\red(\widetilde G)=G$, and $jG\mapsto[j]\widetilde G$ is a
homomorphism (so $\delta_{\tau_{\mathrm{can}}}\equiv 0$). Explicitly
$\widetilde G=[m]\tau(G)$ for any section $\tau$, where $m\equiv 1\ (\bmod\,n)$,
$m\equiv 0\ (\bmod\,p^{e-1})$.
\end{lemma}

\begin{proof}
$E(\Z/p^e)[n]$ injects into $E(\Fp)[n]$ (an $n$-torsion element of
$\Ehat$ is killed by $\gcd(n,p^{e-1})=1$) and surjects onto it (étale
$n$-torsion lifts by Hensel), so reduction is an isomorphism on
$n$-torsion and $\widetilde G$ is the unique order-$n$ lift of $G$.
Writing $\tau(G)=\widetilde G+c$ with $c\in\Ehat$,
$[m]\tau(G)=[m\bmod n]\widetilde G+[m\bmod p^{e-1}]c=\widetilde G$.
\end{proof}

\begin{theorem}[Inertness]\label{thm:main}
Let $\tau$ be any section and $c:=\tau(G)-\widetilde G\in\Ehat$. Then:
\begin{enumerate}
\item[(a)] $\delta_\tau(k)=[k]c-c_k$, where $c_k:=\tau(kG)-[k]\widetilde G
\in\Ehat$ depends only on the point $kG$.
\item[(b)] $[k]c$ depends on $k$ only through $k\bmod p^{e-1}$;
$\log([k]c)=k\log(c)$ has additive order a power of $p$ dividing
$p^{e-1}$.
\item[(c)] Since $\gcd(n,p^{e-1})=1$, the residues $k\bmod n$ (fixed by
$Q$) and $k\bmod p^{e-1}$ (governing $[k]c$) are independent coordinates
of $k$; no value of $[k]c$ is determined by $(E,G,Q)$.
\item[(d)] $\delta_\tau$ is a well-defined function of the instance
(i.e.\ of $k\bmod n$) iff $c=0$ (i.e.\ $\tau$ canonical on
$\langle G\rangle$), in which case $\delta_\tau\equiv 0$. For $c\neq 0$,
$\delta_\tau(k+n)-\delta_\tau(k)=[n]c=[n]\tau(G)\neq O$, the exact
antisymmetry constant of \eqref{eq:antisym}.
\end{enumerate}
\end{theorem}

\begin{proof}
(a) $[k](\widetilde G+c)-([k]\widetilde G+c_k)=[k]c-c_k$. (b) $c\in\Ehat$
has order dividing $p^{e-1}$; apply $\log$. (c) CRT. (d) $c_k$ is a
function of $kG$, hence of $k\bmod n$; thus $\delta_\tau$ factors through
$k\bmod n$ iff $k\mapsto[k]c$ does, iff $[n]c=O$, iff (as
$\gcd(n,p^{e-1})=1$) $c=O$. The shift equals $[n]c=[n]\tau(G)$ because
$[n]\widetilde G=O$.
\end{proof}

\begin{corollary}
For $\gcd(n,p)=1$ no section lift error is a function of $(E,G,Q)$ that
depends on the secret: it is either identically zero (canonical) or
multivalued with ambiguity $[n]\tau(G)\neq O$ (non-canonical). The
$p$-adic canonical lift moves the ECDLP without weakening it.
\end{corollary}

This is exactly why the Gowers anomaly was forced to be artifact-or-%
identity. The only secret-bearing object, $[k]c$, is invisible to the
attacker; the visible $\delta$ is either zero or, taken as a function of
the integer $k$, periodic-with-shift $[n]\tau(G)$---and that shift is the
reflection \eqref{eq:antisym} whose Fourier signature Phase 37
rediscovered.

\section{The boundary is sharp}\label{sec:boundary}

Theorem~\ref{thm:main} hinges on the single inequality $\gcd(n,p)=1$.
Set $n=p$ (an anomalous curve, $\#E(\Fp)=p$) and every clause degenerates
toward the attacker: $k\bmod n=k\bmod p$ now coincides with the residue
the formal group sees, and the obstruction $[n]c=[p]c$ \emph{vanishes}
because $c\in\Ehat(p\Z/p^2)$ has order dividing $p$. The lift error
becomes well-defined and the formal logarithm of $[p]\cdot(\text{lift})$
is a homomorphism $E(\Fp)\to\Fp$ recovering $k$---Smart's attack.

We verified both sides numerically (Table~\ref{tab:boundary}). On every
ordinary test curve the obstruction $[n]\tau(G)$ is a nonzero kernel
element of \emph{maximal} order $p^{e-1}$ (valuation $1$); on an
anomalous curve $y^2=x^3+4x+7$ over $\mathbb F_{53}$ the obstruction
vanishes and Smart's attack recovers $k$ with $100\%$ success ($12/12$
random secrets). The same code path that is inert for $\gcd(n,p)=1$
becomes a total break the instant $n=p$.

\begin{table}[t]
\centering
\small
\begin{tabular}{lcccc}
\toprule
Regime & curves & $\gcd(n,p)$ & obstruction $[n]\tau(G)$ & ECDLP\\
\midrule
Ordinary (ECDSA) & $p=31$--$163$, $j{=}0$ & $1$ &
  nonzero, $\ord=p^{e-1}$ & inert\\
Anomalous (Smart) & $y^2{=}x^3{+}4x{+}7/\mathbb F_{53}$ & $p$ &
  $O$ (vanishes) & broken, $12/12$\\
\bottomrule
\end{tabular}
\caption{Phases 41--42. The well-definedness obstruction of
Theorem~\ref{thm:main}(d) is nonzero of maximal order on ordinary curves
and vanishes on anomalous curves, where Smart's attack succeeds. The
hypothesis $\gcd(n,p)=1$ is the exact boundary.}
\label{tab:boundary}
\end{table}

\section{Numerical verification}\label{sec:verification}

Theorem~\ref{thm:main} was checked on the secp256k1 family
$y^2=x^3+7$ at $e=4$ for ten primes $p=31,\dots,163$ with $\gcd(n,p)=1$
(Phase 41). All kernel arithmetic was done through the additive
$z$-coordinate to avoid projective degeneracy on deep kernel points. On
every prime: the obstruction $C_n:=z([n]\tau(G))$ is nonzero and lies in
$p\Z$ with valuation $1$, of additive order exactly $p^{e-1}$ (clause d,
b); the antisymmetry $z(\delta(k))+z(\delta(n-k))\equiv C_n$ holds for
all $k$ (clause d / \eqref{eq:antisym}); the integer-representative shift
$z(\delta(k+n))-z(\delta(k))\equiv C_n$ holds on all eight
arithmetically-reliable primes (clause d; it is a group identity, and on
the two remaining primes the deep multiply degenerates in exact
arithmetic, flagged by a self-consistency probe); and $\gcd(n,p^{e-1})=1$
(clause c). The canonical-lift case $\delta_{\mathrm{can}}\equiv 0$
(Lemma~\ref{lem:can}) was confirmed independently in an exact $p$-adic
computation. Anomalous recovery is Phase 42 (\S\ref{sec:boundary}).

\section{Discussion}\label{sec:discussion}

\paragraph{For negative cryptanalysis.} The reusable lesson is narrow and
practical. (i) A control tests exactly one alternative---the one its null
recomputes. Before trusting it, enumerate the mechanisms that could
generate the observed statistic and check whether the control separates
them; the permutation test here failed precisely because three
mechanisms collapsed under it. (ii) Build one null per mechanism,
\emph{derived from how that mechanism generates the statistic}: H1 from
the prefix-sum structure (shuffle increments, re-sum), H2 from the
reflection's reach (restrict to a half). (iii) Beware cumulative
statistics: any running sum manufactures low-frequency autocorrelation,
so ``looks autocorrelated'' is the null, not the finding. (iv) Where a
structural theorem is available, it bounds what controls can ever find;
Theorem~\ref{thm:main} says no honest control will turn $\delta$ into an
oracle for $\gcd(n,p)=1$.

\paragraph{For automated and AI-assisted exploration.} Pipelines that
fan out many structural probes and surface the ``most anomalous'' are
increasingly common, and they are false-positive engines: with enough
encodings and transforms, an apparent $5$--$11\sigma$ effect is cheap.
The safeguard is not more probes but mechanism-derived nulls and, where
possible, a closed-form reason the probe must be inert. The value of the
theorem in this study is exactly that: it converts ``we tried forty
things and they failed'' into ``here is the one-line reason they had to.''

\paragraph{Honest scope.} Theorem~\ref{thm:main} is not new mathematics.
It is an explicit instance, in the lift-error language, of the principle
of \cite{silverman4faces,silvermanlift,gadiyarpadma} that lifting
transports the ECDLP; the anomalous degeneration is Smart's attack
\cite{smart}. It concerns section lift errors only; pairing-based
functionals are a separate mechanism (inert here for the independent
reason of huge embedding degree). We claim novelty only for the
methodology and for the explicit identification of the obstruction with
the antisymmetry constant.

\section{Conclusion}

A long negative exploration of $p$-adic ECDLP lifting produced one
seductive anomaly. A standard control endorsed it; mechanism-derived
controls killed it, attributing it entirely to an exact public identity;
and an inertness theorem explains why no such anomaly on ordinary curves
can ever be more than artifact-or-identity, with the boundary falling
exactly at the anomalous curves where the lift attack genuinely works.
The cryptanalytic verdict is the expected negative one. The transferable
result is a discipline: name your alternative hypotheses, build a null
for each from its mechanism, and prefer a structural reason for inertness
over an accumulation of passed tests.

\begin{thebibliography}{9}
\bibitem{silverman4faces} J.~H.~Silverman. \emph{The Four Faces of
Lifting for the Elliptic Curve Discrete Logarithm Problem.} Talk, ECC
2007.
\bibitem{silvermanlift} J.~H.~Silverman. \emph{Lifting and Elliptic Curve
Discrete Logarithms.} SAC 2008, LNCS 5381, Springer, 2009.
\bibitem{gadiyarpadma} H.~G.~Gadiyar and R.~Padma. \emph{The discrete
logarithm problem over prime fields: the safe prime case. The Smart
attack, non-canonical lifts and logarithmic derivatives.}
arXiv:1702.07107, 2017. (See also Gadiyar, Maini, Padma, INDOCRYPT 2004.)
\bibitem{smart} N.~P.~Smart. \emph{The discrete logarithm problem on
elliptic curves of trace one.} J.~Cryptology 12 (1999), 193--196.
\bibitem{satoharaki} T.~Satoh and K.~Araki. \emph{Fermat quotients and
the polynomial time discrete log algorithm for anomalous elliptic
curves.} Comm.\ Math.\ Univ.\ Sancti Pauli 47 (1998), 81--92.
\bibitem{semaev} I.~A.~Semaev. \emph{Evaluation of discrete logarithms on
some elliptic curves.} Math.\ Comp.\ 67 (1998), 353--356.
\bibitem{lst} J.~Lubin, J.-P.~Serre, J.~Tate. \emph{Elliptic curves and
formal groups.} Woods Hole, 1964.
\bibitem{satohcount} T.~Satoh. \emph{The canonical lift of an ordinary
elliptic curve over a finite field and its point counting.} J.~Ramanujan
Math.\ Soc.\ 15 (2000), 247--270.
\bibitem{silvermanaec} J.~H.~Silverman. \emph{The Arithmetic of Elliptic
Curves}, 2nd ed., GTM 106, Springer, 2009.
\bibitem{greentao} W.~T.~Gowers; B.~Green and T.~Tao. \emph{Inverse
theorems for the Gowers norms.} (For the $\Uk^2$/Fourier correspondence
used in \S\ref{sec:falsepos}.)
\end{thebibliography}

\end{document}
